\section{Related Work}
\label{sec:related}
% Intent: position Anchor against four lines of work and state what is not
% new. Every claim about a cited paper must be checked against the paper
% before submission (references.bib marks unverified entries).

\paragraph{Compounding error and multi-step training.}
A model trained for one-step prediction and rolled out feeds its own errors back as inputs, and the error compounds with the horizon. Remedies include training on the model's own rollouts \citep{talvitie2014model, venkatraman2015improving}, multi-step losses against the observed future, as in latent overshooting \citep{hafner2019planet} and in exponentially weighted multi-step objectives that help under observation noise \citep{benechehab2024multistep}, and direct prediction of distant states without recursion \citep{lambert2021learning}. Under Markov misspecification direct prediction is statistically preferable to the rollout, while a well-specified one-step model is the more efficient choice \citep{somalwar_statistical}; our AR(1) analysis (Proposition~\ref{prop:t4-ar1}) gives a closed form of this effect. All of these losses use the true future as the target, i.e.\ the Monte Carlo end of the family of Section~\ref{sec:method}. \todo{Jumpy world models \citep{farebrother2026jumpy}: read and position.}

\paragraph{Self-predictive representations and latent consistency.}
Latent world models used for control are trained to predict their own future latents: SPR \citep{schwarzer2021spr} and TD-MPC \citep{hansen2022tdmpc, hansen2024tdmpc2} compare the multi-step latent rollout with the encoding of the observed future under a stop-gradient or target encoder, following BYOL \citep{grill2020byol}. The theory of these objectives is one-step: \citet{tang2023selfpredictive} show that the stop-gradient prevents collapse and characterize the learned representation spectrally, and \citet{ni2024bridging} show that one-step self-prediction yields latents sufficient for control, including from histories, and leave multi-step prediction without analysis. Our results extend this line to multi-step targets: the residual lemma, harmlessness for Markov latents (Proposition~\ref{prop:t1}) and the strict improvement under a Markov gap (Proposition~\ref{prop:t2-ar1}) hold for the Monte Carlo targets these methods use as well as for bootstrapped ones, and without the stop-gradient a bootstrapped target moves the one-step optimum even for a Markov latent \todo{population fixed points computed; confirm with the e02 ablation D\_nosg}. \todo{Check how SPR and TD-MPC2 weight horizons and whether their targets are encoders with or without history.}

\paragraph{Bootstrapped predictions.}
Learning a prediction from another prediction made at a later time is the idea of temporal-difference networks \citep{sutton2004tdnetworks, tanner2005tdlambda}, which also mix bootstrapped and Monte Carlo targets through eligibility traces, and of predictive state representations and general value functions as models of the world \citep{littman2002psr, sutton2011horde}. Anchor brings this idea to latent world models with a one-step dynamics model: the target is the model's own rollout from a later encoded observation, mixed over anchor points as the $\lambda$-return mixes $n$-step returns. What is new is the analysis of what the mixture does to the rollout of the world model, not the bootstrapped target itself.

\paragraph{Partial observability and adaptive mixing.}
For value functions, the difference between TD($\lambda$) estimates with different $\lambda$ is zero for Markov observations and almost always non-zero otherwise, and minimizing it helps learn memory \citep{allen2024lambda}. Our critic plays the corresponding role for world-model state prediction: it measures the rollout bias caused by the Markov gap without the target variance (Corollary~\ref{cor:critic}). Choosing the trace parameter from bias and variance estimates has been proposed for value prediction \citep{white2016greedy, mann2016adaptive}; Anchor chooses the weights of its anchor points in the same spirit, with the bias estimated by the difference of two critics. \todo{Verify how \citet{white2016greedy} and \citet{mann2016adaptive} estimate bias and variance, and state the difference precisely.}

\paragraph{What is not new.}
Compounding error, direct multi-step prediction, horizon-conditioned predictors, multi-step consistency losses with the true future, stop-gradient targets and bootstrapped predictions all predate this work. Our contribution is their analysis for latent world models (which targets help, when, and why the stop-gradient is needed), the critic as a measure of the Markov gap, and the bias--variance mixture over anchor points that follows from it.
